\section{Results}
\label{sec:results}

\begin{lede}
Five suites, five seeds, 400 timesteps, four maps, five planners. This section
puts the verdict first --- one table that says, for every rival on every map,
whether SPAR-PIBT delivers more work or less --- and only then explains where
each number comes from.
\end{lede}

\subsection{The answer, before the argument}
\label{sec:results-answer}

\begin{verdict}[Does SPAR-PIBT beat the alternatives?]
\begin{enumerate}[leftmargin=1.3em]
\item \textbf{Against the two published lifelong baselines: yes, everywhere,
  by 13$\times$ to 315$\times$}, at a twentieth to a three-hundredth of their
  cost per timestep, at the strongest significance five seeds can express
  ($p = 0.008$). Token Passing delivers \emph{nothing at all} on
  \variant{corridors}. So does plain PIBT beat them, and that gap belongs to
  PIBT, not to this work.
\item \textbf{Against plain lifelong PIBT --- the planner SPAR extends --- the
  answer depends on the map, and the split is the contribution.} SPAR's aisle
  layer delivers \textbf{27\% and 21\% more} on the two aisle-constrained maps,
  and \textbf{18\% and 19\% less} on the two open-floor maps. At five seeds the
  two losses are significant and \emph{the two wins are not yet}
  ($p = 0.056$ and $0.063$).
\item \textbf{Inside SPAR, the design decision the method is named for is the
  largest effect measured anywhere in this system}: pricing counterflow rather
  than forbidding it is worth $1.9\times$ to $3.7\times$ throughput, at
  $p = 0.008$, on every map where an aisle ever commits a direction.
\end{enumerate}
\end{verdict}

\begin{figure}[H]
\centering
\includegraphics[width=\linewidth]{scorecard.pdf}
\caption{One cell per rival per map: SPAR-PIBT's throughput divided by that
rival's. Blue means SPAR delivers more, red means it delivers less, and the
depth of the colour is the size of the gap. This figure is the whole result;
everything after it is the explanation.}
\label{fig:scorecard}
\end{figure}

\begin{table}[H]
\centering\footnotesize
\begin{tabular}{@{}llrrrrl@{}}
\toprule
 & & \multicolumn{4}{c}{\textbf{Tasks delivered per 1000 timesteps}} & \\
\cmidrule(lr){3-6}
\textbf{Map} & \textbf{Class} & \textbf{SPAR-PIBT} & \textbf{plain PIBT} &
\textbf{TP} & \textbf{RHCR} & \textbf{Verdict for SPAR} \\
\midrule
bottleneck & aisle-constrained & \textbf{149} & 118 & 8 & 12 &
  \wins\ beats all four rivals \\
corridors & aisle-constrained & \textbf{158} & 130 & 0 & 0.5 &
  \wins\ beats all four rivals \\
narrow & open floor & 289 & \textbf{354} & --- & --- &
  \loses\ behind plain PIBT \\
medium & open floor & 405 & \textbf{502} & 27 & 14 &
  \loses\ behind plain PIBT \\
\bottomrule
\end{tabular}
\caption{The whole evaluation in one table. Throughput is stated per 1000
timesteps rather than per timestep, so 149 means 149 tasks delivered out of
every 1000 steps of simulated time. \textbf{TP} is the better of the two Token
Passing variants on that map. A dash means that baseline suite was not run on
that map. SPAR-PIBT here is the \variant{aisle\_direction\_only} configuration,
which is the best-throughput SPAR configuration on all four maps.}
\label{tab:scorecard}
\end{table}

\begin{caveat}
\textbf{How to read every number in this section.}
\begin{itemize}[leftmargin=1.2em]
\item \textbf{Throughput is tasks per 1000 timesteps} throughout this section
  and every figure. The experiment records tasks per timestep; multiplying by
  1000 changes nothing but the number of leading zeros.
\item \textbf{Higher is better} for throughput, and it is the metric that
  decides. Service time is \emph{censored} and cannot rank planners on its own
  (\Cref{fig:censoring}); cost is milliseconds per timestep, where lower is
  better.
\item \textbf{$p$ is a two-sided exact permutation test} over the same five
  seeds. Five seeds a side gives $\binom{10}{5} = 252$ label splits, so
  \textbf{the smallest attainable $p$-value is $2/252 \approx 0.008$}: a cell
  reading $p = 0.008$ is as significant as this design can report, and a
  genuine difference of moderate size can easily land at $p \approx 0.06$.
\item \textbf{Bold marks a difference at $p \leq 0.05$}; \wins\ means SPAR
  delivers more, \loses\ that it delivers less, \unclear\ that the arms are
  indistinguishable or the comparison did not run.
\item \textbf{``SPAR-PIBT'' means one fixed configuration},
  \variant{aisle\_direction\_only}, unless a table says otherwise. It is the
  best-throughput SPAR configuration on every map in this dataset, so no
  comparison in this section is cherry-picked per map.
\end{itemize}
\end{caveat}

\subsection{Setup}

Everything below comes from one dataset, \code{docs/data/}, written by
\code{experiments/run\_all.py} and committed alongside this document, so every
figure regenerates from the same runs.

\begin{table}[H]
\centering\footnotesize
\begin{tabular}{@{}lllrrl@{}}
\toprule
\textbf{Map} & \textbf{Class} & \textbf{Shape} & \textbf{Robots} &
\textbf{$\lambda$} & \textbf{Served} \\
\midrule
\variant{warehouse\_bottleneck} & aisle-constrained & two halves, one joining corridor & 16 & 0.8 & 18\% \\
\variant{warehouse\_corridors} & aisle-constrained & five parallel single-file corridors & 35 & 1.0 & 13\% \\
\variant{warehouse\_narrow} & open floor & four aisle banks, three cross-aisles & 30 & 1.2 & --- \\
\variant{warehouse\_medium} & open floor & open grid, four cross-aisles & 40 & 1.5 & 33\% \\
\bottomrule
\end{tabular}
\caption{The scenarios, and the classification used throughout this section.
\emph{Aisle-constrained} means a robot in a contended aisle has no way around
it; \emph{open floor} means cross-aisles give it one. ``Served'' is the best
planner's completed tasks over the tasks the arrival process released --- the
saturation of \Cref{sec:regime}, measured.}
\label{tab:scenarios}
\end{table}

Five seeds per cell, Poisson arrivals, and the statistics of
\Cref{sec:statistics} throughout.

\subsection{Against the published baselines}
\label{sec:results-baselines}

\begin{figure}[H]
\centering
\includegraphics[width=\linewidth]{headline.pdf}
\caption{Throughput on the three maps the baseline suite covers. Both
baselines are the two short bars on every panel.}
\label{fig:headline}
\end{figure}

\begin{table}[H]
\centering\small
\begin{tabular}{@{}llrrrrl@{}}
\toprule
\textbf{Map} & \textbf{Planner} & \textbf{$\Theta_{1000}$} & \textbf{$\times$ SPAR} &
\textbf{svc} & \textbf{ms/step} & \textbf{Verdict} \\
\midrule
\multirow{5}{*}{bottleneck}
 & \textbf{SPAR-PIBT} & \textbf{149} & --- & 157 & 0.61 & --- \\
 & plain lifelong PIBT & 118 & $1.27\times$ & 154 & 0.41 & \unclear\ $p = 0.056$ \\
 & Token Passing & 8 & $19.9\times$ & 31 & 36.2 & \wins\ $p = 0.008$ \\
 & Token Passing + recovery & 8 & $17.5\times$ & 42 & 32.3 & \wins\ $p = 0.008$ \\
 & RHCR & 12 & $13.0\times$ & 23 & 75.0 & \wins\ $p = 0.008$ \\
\midrule
\multirow{5}{*}{corridors}
 & \textbf{SPAR-PIBT} & \textbf{158} & --- & 163 & 1.38 & --- \\
 & plain lifelong PIBT & 130 & $1.21\times$ & 157 & 0.82 & \unclear\ $p = 0.063$ \\
 & Token Passing & 0 & $\infty$ & --- & 80.8 & \wins\ $p = 0.008$ \\
 & Token Passing + recovery & 0 & $\infty$ & --- & 40.3 & \wins\ $p = 0.008$ \\
 & RHCR & 0.5 & $315\times$ & 9 & 43.0 & \wins\ $p = 0.008$ \\
\midrule
\multirow{5}{*}{medium}
 & \textbf{SPAR-PIBT} & 405 & --- & 158 & 1.95 & --- \\
 & plain lifelong PIBT & \textbf{502} & $0.81\times$ & 139 & 1.01 & \loses\ $p = 0.008$ \\
 & Token Passing & 4 & $101\times$ & 50 & 222.2 & \wins\ $p = 0.008$ \\
 & Token Passing + recovery & 27 & $15.0\times$ & 145 & 191.2 & \wins\ $p = 0.008$ \\
 & RHCR & 14 & $27.9\times$ & 39 & 97.9 & \wins\ $p = 0.008$ \\
\bottomrule
\end{tabular}
\caption{Five seeds, 400 timesteps. $\Theta_{1000}$ is tasks delivered per 1000
timesteps; ``$\times$ SPAR'' is SPAR-PIBT's throughput divided by that row's, so
$19.9\times$ means SPAR delivers twenty times the work; ``svc'' is mean service
time over completed tasks only (\Cref{sec:metrics}); $p$ is a two-sided
permutation test on throughput against SPAR-PIBT. A dash in the service-time
column means the planner finished no tasks at all, so the statistic does not
exist.}
\label{tab:baselines}
\end{table}

Three things in that table are worth reading carefully.

\paragraph{Both baselines lose by one to two orders of magnitude, significantly, everywhere.}
Every baseline row sits at $p = 0.008$, the floor of a five-seed test. Token
Passing completes \emph{literally zero} tasks on \variant{warehouse\_corridors}
across all five seeds. Its counters (\code{tp\_path\_not\_found},
\code{tp\_forced\_holds}) confirm genuine gridlock rather than a crash: robots
queue nose-to-tail and no robot can reserve a path through the robots ahead of
it, for the rest of the run. This is the structural failure of
\Cref{sec:related-tp}, and the first animation in \code{docs/gifs/} shows it
beside the same scenario under priority inheritance.

\paragraph{The service-time column is a trap, and this table is where it springs.}
RHCR reports a mean service time of $9$ timesteps on
\variant{warehouse\_corridors} --- easily the best figure in the table --- while
delivering half a task per 1000 timesteps. It is not fast; it is selective, and
only by accident. Token Passing's dashes are the same phenomenon taken to its
limit. \textbf{Never rank two planners by service time unless their throughputs
are comparable.}

\begin{figure}[H]
\centering
\includegraphics[width=\linewidth]{censoring.pdf}
\caption{The same five planners, twice. Ranking by the right panel alone
inverts the ranking that matters.}
\label{fig:censoring}
\end{figure}

\paragraph{Cost does not buy throughput here.}
SPAR's aisle layer costs 0.6--2.0\,ms per timestep against plain PIBT's
0.4--1.0 (4.3\,ms with every mechanism enabled); both baselines cost
32--222\,ms and deliver almost nothing at this scale.
\Cref{sec:results-caveats} states what that does and does not say about RHCR.

\begin{figure}[H]
\centering
\includegraphics[width=\linewidth]{cost-benefit.pdf}
\caption{Planner cost against what it delivers; up and to the left is better.
Note the log scale on cost.}
\label{fig:cost-benefit}
\end{figure}

\subsection{Against the planner SPAR extends}
\label{sec:results-honest}

The baselines are not the hard comparison. Plain lifelong PIBT is, because
SPAR's aisle layer is built on top of it and has to earn its place there.

\begin{table}[H]
\centering\footnotesize
\begin{tabular}{@{}llrrrll@{}}
\toprule
\textbf{Map} & \textbf{Class} & \textbf{SPAR} & \textbf{plain PIBT} &
\textbf{difference} & \textbf{$p$} & \textbf{Verdict} \\
\midrule
bottleneck & aisle-constrained & \textbf{149} & 118 & $+27\%$ & 0.056 &
  \unclear\ ahead, not yet significant \\
corridors & aisle-constrained & \textbf{158} & 130 & $+21\%$ & 0.063 &
  \unclear\ ahead, not yet significant \\
narrow & open floor & 289 & \textbf{354} & $-18\%$ & 0.016 &
  \loses\ behind, significant \\
medium & open floor & 405 & \textbf{502} & $-19\%$ & 0.008 &
  \loses\ behind, significant \\
\bottomrule
\end{tabular}
\caption{The best SPAR configuration on each map --- \variant{aisle\_direction\_only}
in all four --- against plain lifelong PIBT, in tasks per 1000 timesteps. The
pattern is consistent with the design; \emph{the evidence for the winning half
of it is currently thin}, and \Cref{fig:winloss} shows the same split across
every other configuration.}
\label{tab:honest}
\end{table}

Read as evidence rather than as a headline, this table says four things.

\begin{enumerate}
\item \textbf{The split is by map class, not by map}, and it is the split the
  method predicts: the aisle layer helps where a contended aisle has no
  alternative route around it, and costs where cross-aisles already provide one.
  Four maps is four points, and \Cref{sec:results-caveats} says so.
\item \textbf{The two losses are significant; the two wins are not}
  ($p = 0.056$ and $0.063$ against $0.016$ and $0.008$). At five seeds a $+27\%$
  difference simply cannot reach the floor unless the seeds barely overlap.
  Ten seeds would settle it either way, and that is the obvious next run rather
  than a caveat to argue around.
\item \textbf{Throughput is not the only thing the aisle layer moves.} On both
  maps where it wins, it also drives about a quarter less to deliver the same
  task and halves the head-on conflicts per delivery
  (\Cref{sec:results-other-metrics}) --- effects that are significant where the
  throughput difference is not.
\item \textbf{The cost of the win is legible too}: backtracks per delivered task
  rise 85--130\% wherever an aisle commits a direction. Pricing counterflow makes
  the inheritance recursion work harder, and that is the method's largest cost.
\end{enumerate}

\begin{figure}[H]
\centering
\includegraphics[width=\linewidth]{winloss.pdf}
\caption{Every configuration against plain lifelong PIBT, with a permutation
test on the same seeds. The two left columns are the aisle-constrained maps.}
\label{fig:winloss}
\end{figure}

Two rows of \Cref{fig:winloss} deserve comment. \variant{turning\_cost\_only}
--- a flat anti-zigzag penalty with nothing to do with aisles --- is the best
variant on \variant{warehouse\_corridors} (196 per 1000, $p = 0.008$) and the
second best on \variant{warehouse\_medium}. The cheapest mechanism in the system
is among the most effective, and any honest account of this work has to say so.
And on \variant{warehouse\_bottleneck} every variant beats plain PIBT by 21--28\%
while none of them clears $p = 0.05$: with no aisle ever committing a direction,
those variants differ from the base only in the turning cost, and five seeds
cannot separate them.

\subsection{Saturation, in one picture}

\begin{figure}[H]
\centering
\includegraphics[width=\linewidth]{load.pdf}
\caption{Offered demand against served demand. Between 13\% and 33\% of the
tasks released are delivered; the rest are still queued when the horizon ends.}
\label{fig:load}
\end{figure}

This is \Cref{sec:regime} measured rather than asserted, and it is the context
every other number needs. Throughput here is capacity, not activity.

\subsection{Which mechanisms earn their cost}
\label{sec:results-mechanisms}

\begin{figure}[H]
\centering
\includegraphics[width=\linewidth]{forest.pdf}
\caption{Each of the four repaired mechanisms against the same configuration
without it, on every map. Blue is $p < 0.05$.}
\label{fig:forest}
\end{figure}

\begin{table}[H]
\centering\footnotesize
\begin{tabular}{@{}lrrrr@{}}
\toprule
\textbf{Mechanism (with / without)} & \textbf{bottleneck} & \textbf{corridors} & \textbf{narrow} & \textbf{medium} \\
\midrule
direction \textbf{ranks} vs \textbf{rejects} & 149 / 149 & \textbf{158 / 84} & \textbf{289 / 78} & \textbf{405 / 161} \\
same, with reservations & \textbf{149 / 60} & \textbf{95 / 43} & \textbf{152 / 33} & 217 / 88 \\
bounded maximum green & 149 / 149 & 158 / 159 & 289 / 282 & 405 / 360 \\
recovery needs corroboration & 149 / 148 & \textbf{134 / 22} & 204 / 134 & 245 / 177 \\
\midrule
\multicolumn{5}{@{}l@{}}{\emph{as a ratio, mechanism on over mechanism off}} \\
direction \textbf{ranks} vs \textbf{rejects} & $1.0\times$ & $\mathbf{1.9\times}$ & $\mathbf{3.7\times}$ & $\mathbf{2.5\times}$ \\
same, with reservations & $\mathbf{2.5\times}$ & $\mathbf{2.2\times}$ & $\mathbf{4.6\times}$ & $2.5\times$ \\
bounded maximum green & $1.0\times$ & $1.0\times$ & $1.0\times$ & $1.1\times$ \\
recovery needs corroboration & $1.0\times$ & $\mathbf{6.1\times}$ & $1.5\times$ & $1.4\times$ \\
\bottomrule
\end{tabular}
\caption{Throughput per 1000 timesteps, with the mechanism / without it, five
seeds. \textbf{Bold} is $p \leq 0.05$. The lower block is the same rows as a
ratio, which is the form the argument uses.
\code{config.PAIRED\_DESIGNS}.}
\label{tab:mechanisms}
\end{table}

\paragraph{Direction as a price rather than a rule is the largest effect in the system.}
On the three maps where an aisle ever commits a direction it is worth
$1.9\times$ (corridors), $2.5\times$ (medium) and $3.7\times$ (narrow)
throughput, each at $p = 0.008$; with entry admission also on, the range runs
to $4.6\times$. On \variant{warehouse\_bottleneck} the two arms are identical to
three decimals, and that is not a null result: 16 robots spread over 25 aisles
never push the demand imbalance past $\tau_{\mathrm{switch}} = 5.0$, so no aisle
ever commits a direction and there is nothing for the treatment to change.
\textbf{That cell is an experiment that did not run}, and is reported as such
rather than as evidence of no effect.

This is \Cref{sec:contribution-soft} measured. The second animation in
\code{docs/gifs/} is the corridors pair, side by side.

\paragraph{Corroborated stall detection is the second.}
On \variant{warehouse\_corridors}, 134 against 22 tasks per 1000 timesteps
($p = 0.008$): a six-fold difference obtained entirely by refusing to escalate
on the weakest of the three stall signals. On the other maps the effect points
the same way and does not clear five-seed significance ($p = 0.056$ on narrow,
$0.230$ on medium). The mechanism matters most exactly where deadlocks are
easiest to manufacture.

\paragraph{The maximum green is not significant on throughput, and that is the finding.}
Its throughput effect is between $-1$ and $+45$ per 1000 timesteps, with $p$
from $0.437$ to $0.968$. It was decisive while direction was a hard constraint,
because a starved robot then had no way out at all; once counterflow is merely
expensive, a robot can buy its way out of a starved aisle and liveness no longer
rests on the signal. The two fixes are partly redundant. What the maximum green
does change is the signal itself: on \variant{warehouse\_narrow} it produces
$22.5$ direction switches per 1000 steps against $1.0$ without it
($p = 0.016$), of which $5.2$ per run are flips forced by the rule rather than
by demand (\code{starvation\_flips}). Without it, those aisles simply never turn.

\subsection{The ablation ladder is not monotone}

\begin{figure}[H]
\centering
\includegraphics[width=\linewidth]{ablation-ladder.pdf}
\caption{The cumulative ladder of \Cref{sec:ablation-ladder}. Adding
mechanisms does not monotonically help, and on two maps the first rung is the
best one.}
\label{fig:ladder}
\end{figure}

This is why \Cref{tab:scorecard} reports one named configuration rather than
``the full method'': stacking every mechanism is not the best version of this
system on any map in this dataset, and saying otherwise would misrepresent it.

\subsection{Hypotheses}
\label{sec:results-hypotheses}

\begin{figure}[H]
\centering
\includegraphics[width=\linewidth]{hypotheses.pdf}
\caption{Each hypothesis against the control that isolates its mechanism, on
the quantity it actually claims to move.}
\label{fig:hypotheses}
\end{figure}

Scoring each hypothesis on its own metric rather than on global throughput is
what makes these verdicts mean anything: H3 claims \emph{fewer head-on
conflicts} and H1 claims \emph{narrow aisles flow better}, and judging either by
throughput can miss a mechanism doing exactly what it promises, or credit one
that is not.

\begin{itemize}
\item \textbf{H2 (hysteresis reduces switching) is supported on every map}, at
  the five-seed floor, and by two orders of magnitude: 14.5 switches per 1000
  steps against 2188 on \variant{warehouse\_medium}.
\item \textbf{H3 (entry admission reduces head-on conflicts) is map-dependent}:
  supported on corridors (702 against 1130), contradicted on narrow. Both
  directions are real and the map decides which.
\item \textbf{H1 is contradicted on its own metric on corridors} and flat
  elsewhere, while the same configuration raises \emph{global} throughput on
  three of four maps. Both are true and not in tension: one-way aisles route
  traffic through fewer aisle transits per delivered task, so per-aisle flow
  falls while end-to-end flow rises. The lesson is about the denominator, not
  the mechanism.
\item \textbf{H5's mechanism is inert.} Changing $R_{\mathrm{near}}$ and
  $R_{\mathrm{far}}$ produces runs identical to the last decimal, because
  $\beta$ only breaks ties among candidates with \emph{equal} progress and the
  decay lowers it precisely where $\alpha \Delta$ already decides. The weight
  is not inert, and is net harmful: \variant{no\_direction\_term} ($\beta = 0$)
  beats \variant{aisle\_direction\_only} on medium (433 against 405 per 1000)
  and ties it on narrow.
\item \textbf{H4 and H6 are unsupported on three maps each}, and H6 is supported
  on corridors.
\end{itemize}

\subsection{De-confounded findings}
\label{sec:results-factorial}

The $2 \times 2$ decomposition of \Cref{eq:factorial}, on the design that could
not be read at all before the reservation layer's gating was fixed
(\Cref{sec:aisle-reservations}):

\begin{table}[H]
\centering\small
\begin{tabular}{@{}lrrrrr@{}}
\toprule
\textbf{Map} & \textbf{base} & \textbf{aisle direction} & \textbf{reservations} & \textbf{both} & \textbf{interaction} \\
\midrule
corridors & 148 & \textbf{158} & 125 & 95 & $-40$ \\
medium & 337 & \textbf{405} & 262 & 217 & $-112$ \\
\bottomrule
\end{tabular}
\caption{\code{aisle\_direction\_vs\_reservations}, tasks per 1000 timesteps.
The attribution reverses once the cells are distinct runs.}
\label{tab:factorial}
\end{table}

Aisle-level direction control is the flag that \emph{helps} in isolation
($+7\%$ and $+20\%$); entry admission is the one that costs ($-16\%$ and
$-22\%$); and the two interact strongly negatively, so the combined rung of the
ladder cannot be read as one contribution plus the other. The earlier verdict
--- direction blamed, reservations exonerated --- was reading a cell in which
factor B was a bit-identical copy of the base.

\subsection{What the other metrics say}
\label{sec:results-other-metrics}

\code{docs/metrics.md} is the companion to this section: every reported
quantity, which direction is good, and how each one misleads on its own. Three
results there are worth carrying back into this document, because they are the
clearest statement of what the aisle layer buys and what it costs, and none of
them is visible in a throughput column:

\begin{itemize}
\item \textbf{Travel distance per delivered task} falls 23\% on
  \variant{warehouse\_bottleneck} ($p = 0.008$) and 25\% on
  \variant{warehouse\_corridors} ($p = 0.016$) --- the fleet drives a quarter
  less to deliver the same task --- and rises 13\% and 21\% on the two open
  maps. The raw total is lower everywhere and means nothing: a planner that
  delivers less also drives less.
\item \textbf{Head-on conflicts per 1000 steps} fall 34\% on
  \variant{warehouse\_bottleneck} ($p = 0.016$), and 52\% per delivered task.
  That is the mechanism measured on its own terms rather than through
  throughput.
\item \textbf{Backtracks per delivered task} rise by 85--130\% wherever an
  aisle commits a direction ($p = 0.008$). Pricing counterflow makes the
  inheritance recursion explore further before it resolves, and that is the
  largest single cost of the method.
\end{itemize}

Two of these are significant on exactly the maps where the throughput
difference is not, which is the strongest form the winning claim currently
takes: \emph{on aisle-constrained maps the aisle layer demonstrably changes the
traffic in the direction it was designed to, and the throughput consequence is
consistent but not yet significant at five seeds.}

\subsection{Threats to validity}
\label{sec:results-caveats}

\begin{enumerate}
\item \textbf{Task assignment is held constant} across all five planners
  (\Cref{sec:assignment}). This isolates the movement layer, and it also means
  neither baseline was tested with the assignment policy its own paper pairs it
  with.
\item \textbf{RHCR's window and replan period (10/5) are untuned defaults}, not
  values chosen for this scale. RHCR targets hundreds of robots on much larger
  warehouses, where a rolling horizon pays for itself against full CBS. Its poor
  showing here is a property of \emph{this baseline at this scale with these
  defaults}, and is not evidence about RHCR in general.
\item \textbf{Token Passing's gridlock, by contrast, is structural} --- no
  priority inheritance, no backtracking --- and is well documented in the
  lifelong-MAPD literature. It will not go away with different hyperparameters.
\item \textbf{Five seeds is thin}, as \Cref{tab:honest} shows directly. Every
  claim of significance here rests on a test whose floor is $p = 0.008$, and
  every claim of non-significance may simply be underpowered.
\item \textbf{Four maps, one arrival process, one fleet size per map.} The
  boundary this document draws between ``aisle-constrained'' and ``open floor''
  is drawn through four points.
\end{enumerate}
